% Member 2 - final report, III.D Training Protocol
\subsection{Training Protocol}
\label{sec:training_protocol}

Every group is trained with the same protocol, so that the differences in
the results are attributable to the training data and to nothing else.

\paragraph{Architecture.}
The classifier is a ResNet-18 with the standard four stages, initialised
from the ImageNet-pretrained weights. No group uses a different
architecture, except the MixStyle baseline, which adds the published
MixStyle feature-mixing layer; that baseline is trained on the original
images and carries a separate configuration fingerprint, so that it can
never be conflated with a styled group.

\paragraph{Optimisation.}
We use SGD with momentum 0.9, weight decay $5\times10^{-4}$, learning rate
0.01 and a cosine annealing schedule, with label smoothing 0.1 and a batch
size of 64. The step budget is fixed at 660 optimiser steps for every
group. This number is 30 nominal epochs over the 1{,}349-image Original
training set, and it is deliberately a \emph{step} budget rather than an
epoch budget: the styled groups have twice as many training rows, so an
epoch budget would give them twice the number of gradient updates and the
comparison would be confounded.

\paragraph{Data augmentation during training.}
Training images are resized to $224\times224$, randomly cropped with
rejection of degenerate crops, and randomly horizontally flipped, with the
usual ImageNet normalisation. These are the standard training-time
augmentations and they are applied identically to every group.

\paragraph{Which checkpoint is reported.}
We report the \emph{final} checkpoint of the fixed step budget, not the
best one on the validation set. The reason is quantitative and worth
stating: the Photo validation accuracy saturates at 97--98\% within the
first 60 steps, and with 150 validation images one image is worth
$0.67\%$, so selecting the best validation checkpoint selects noise. On our
baseline this selection more than doubled the seed-to-seed spread of the
cross-domain accuracy, from a standard deviation of 2.71 to 7.22 points.
Using the final checkpoint introduces no selection bias at all, and it makes
the protocol identical for every group without exception.

\paragraph{Seeds.}
Every reported group is trained with five seeds, 0 to 4. The seed controls
the weight initialisation, the data order, the crops and the flips. All
seeds of all groups use the same set of seed values, so that every
comparison in this paper is a \emph{paired} one over the same five
realisations; every difference we report is a paired difference, with its
bootstrap confidence interval, its paired $t$-test, a
Benjamini--Hochberg corrected $p$-value across the family of comparisons,
and Cohen's $d$.

\paragraph{Reproducibility.}
Every run records a configuration fingerprint that covers the training
protocol, the evaluation splits and the manifest it used. A run whose
recorded fingerprint differs from the current one is retrained rather than
reused, so a table can never mix two protocols. The fingerprints are
computed from the configuration itself rather than being entered by hand,
which is what makes the check meaningful.
